\section{Conclusion}

This thesis set out to determine whether bounded Android permission evidence could support a useful privacy-review product without being misrepresented as automated legal judgement. The conclusion answers the four research questions, integrates the architectural and methodological contributions, states the final product position, and identifies the next evidence required. Operational deployment details already analysed in Chapter 6 are not repeated here.

\subsection{Answers to the Research Questions}

\textbf{RQ1: deterministic and safe transformation.} Permission summaries are accepted through a conservative evidence contract, rejected when identifiers, numbers, sources, windows, timestamps, or contexts are unsafe, and evaluated through a deterministic rule-pack engine. Temporal isolation addresses replay and overlap within the limits of aggregate windows.

The answer is supported by more than nominal examples. Runtime validation treats TypeScript declarations as developer guidance rather than as a security boundary. Permission identifiers must be owned values from the admitted set; numeric fields must be safe integers; windows must be ordered, bounded, and plausibly timed; and source values must match the closed provenance vocabulary. A rejected observation does not become a normal finding. This fail-closed behaviour is important because silent normalisation would create reassurance from evidence whose meaning is unknown.

Determinism also includes state. An identical record does not accumulate repeatedly, a late record cannot move the repository's retention watermark backwards, and cross-window evaluation uses completed non-overlapping history. These controls make repeated test runs explainable and prevent several aggregate replay errors. The qualification remains that aggregates do not contain stable identities for their constituent events. The implementation can refuse unsafe combination; it cannot reconstruct information the producer did not supply.

\textbf{RQ2: calibrated communication.} Each finding preserves provenance, temporal scope, rule-pack identity, rationale, and missing-context questions. The interface distinguishes needs-review, evidence-gap, and no-technical-concern states, and repeatedly states that technical output is not a legal verdict. Synthetic demonstrations remain labelled throughout the data path.

Calibrated communication is implemented as an information architecture, not as a disclaimer appended to an otherwise absolute score. Overview explains the product's current evidence and offers a clear next action. Findings organises review states and exposes detail progressively. Settings identifies the selected region and provides data controls. Within a finding, status, source, time, numerical reason, pack identity, and follow-up questions appear in an order intended to support both rapid triage and deeper inspection.

The communication result remains an expert assessment of the implemented design. The device render and interaction checks show that the path exists and is visually coherent on the tested handset. They do not prove that a population will understand the distinction between technical concern and legal compliance. A task-based participant study is required to establish that claim. The thesis therefore answers RQ2 at the level of traceable interface design and bounded device evidence, not validated human comprehension.

\textbf{RQ3: regional extensibility.} Jurisdiction-dependent material is loaded from registered, versioned packs behind a stable contract. The EU GDPR pack is the evaluated legal objective; the Research Baseline demonstrates software switching without claiming another jurisdiction. New regional packs require their own legal and empirical validation.

This result separates two forms of extensibility that are often conflated. Software extensibility means that another admitted pack can supply identifiers, labels, thresholds, legal references, rationales, and questions without rewriting navigation, persistence, or Android integration. Substantive legal extensibility means that those contents are correct for a jurisdiction and use context. The project demonstrates the first and defines governance requirements for the second. It does not use a placeholder pack as evidence of comparative-law validity.

The contract also improves maintenance of the evaluated GDPR objective. A finding stores the pack and rule version that produced it, so a later rule update need not erase historical meaning. Registry loading validates mandatory authored, reviewed, effective-date, review-authority, review-scope, release-scope, locale, trigger, pack-lifecycle, source-status, source-version, and source-lifecycle fields; incompatible status--lifecycle combinations, missing versions, impossible dates, guidance-only legal packs, and unknown identifiers fail closed. Basis-specific evidence questions prevent a lawful-basis label alone from producing an evidence-complete result. This design provides a foundation for future jurisdictions, but every production pack still needs independent qualified review, versioned fixtures, localisation, source monitoring, and an update policy.

\textbf{RQ4: supported claims.} Compilation, lint, deterministic and adversarial suites, release APK/AAB assembly, merged-manifest inspection, physical installation, cold launch, and rendered workflows passed for the recorded candidate. These results support an initial store-submission engineering candidate. They do not prove legal compliance, unrestricted Android visibility, long-run reliability, broad usability, accessibility conformance, production signing, or store acceptance.

The answer to RQ4 is therefore a claim map rather than a maturity slogan. Logical suites support engine conformance to specified rules. Package inspection supports statements about identifiers, SDK targets, permissions, backup configuration, and incorporated resources. Physical checks support only the recorded handset, build, and interaction interval. Visual review supports an expert judgement that the candidate is coherent enough for controlled submission work. External owner actions and broader empirical studies remain explicit gates.

This structure makes negative results useful. If a future OEM test misses observations, the acquisition claim can be narrowed without invalidating the regulation-pack architecture. If participants misunderstand a label, the presentation layer can change without rewriting historical evidence. If legal review changes a mapping, a new pack version can be admitted without silently rewriting old findings. The research questions are therefore connected by contracts while retaining separable evidence.

\subsection{Contributions and Research Significance}

The first contribution, the evidence-bounded architecture, establishes the governing principle: every conclusion must retain the source, temporal scope, rule identity, and known missing context needed to interpret it. This principle influences data types, validation, storage, interface language, testing, and release documentation. It is not a wrapper around the implementation. Without it, individual technical features could still work while the product as a whole overstated its authority.

The second contribution, the fail-closed decision pipeline, operationalises that principle at the point where data enters the system. A permissive pipeline would transform malformed or inherited values into ordinary results, creating an especially harmful false sense of safety. Closed identifiers, safe numbers, bounded windows, explicit sources, replay semantics, and structured evidence gaps ensure that uncertainty remains visible. The pipeline is intentionally deterministic so an auditor can reconstruct why a record was accepted and why a finding changed.

The third contribution, the regulation-pack contract, locates jurisdictional interpretation in a reviewable component. This reduces coupling and clarifies the boundary between generic signal processing and legally motivated questions. The contract includes identity and versioning because a rule without provenance cannot support historical accountability. The contribution is the architecture and governance model, not a claim that arbitrary legal systems can be translated automatically.

The fourth contribution, the local-first interface, turns the internal evidence model into a product experience. The three-destination structure removes experimental navigation and emphasises current state, actionable review, and controls. Status is not encoded by colour alone. Synthetic provenance and claim limits are presented near relevant decisions. Local storage, no account requirement, no sensitive runtime permission, disabled backup, and a clear-data action align the product behaviour with its privacy message.

The fifth contribution, the verification record, connects artefacts to claims. Deterministic fixtures, adversarial boundary cases, build outputs, manifest facts, device screenshots, and store-readiness gaps are not collapsed into one test total. This makes the evaluation auditable and limits accidental overstatement. It also defines a repeatable path for future releases: every material change should identify which evidence layers must be renewed.

Methodologically, these contributions form a design-science result: the artefact embodies a solution principle, demonstration shows that the principle can be realised, and evaluation tests the artefact against explicit objectives \cite{hevner2004,peffers2007}. The resulting knowledge claim is conditional rather than universal. Within the recorded evidence sources and Android release candidate, preserving provenance and missing context makes the warning path more auditable. Whether that architecture improves comprehension, decisions, or organisational accountability remains an empirical question for the next study.

Together, the contributions support a form of trustworthy incompleteness. Privacy Lens does not solve every acquisition, legal, or usability problem, but it makes those unsolved parts structurally difficult to hide. In an accountability tool, that property is itself useful. A system that clearly refuses an unsupported verdict can better support responsible human review than a more ambitious classifier whose authority is unclear.

\subsection{Final Product Position}

Privacy Lens 1.2.0 previously received an internal 95.1/100 expert-proxy score under a release-engineering rubric. The stricter publication-target rubric resets the App--thesis predecessor to 70/100, assesses Revision 15 at 79/100, Revision 16 at 84/100, Revision 17 at 87/100, Revision 18 at 89/100, and Revision 19 at 90/100. The difference is intentional: the strict rubric reserves substantial credit for independent legal review, participant and assistive-technology evidence, and publication-quality external replication. Version 1.7.0 uses version code 8 and adds an exact source-version and lifecycle contract to the existing legal-evidence and human-agency architecture. It builds as ARM64 APK/AAB, passes the recorded source and release contracts, renders the governed source register on the OPPO device, and completes five cold-start rounds without a package crash-buffer entry. None of these scores is a university mark, publication decision, legal opinion, accessibility certification, or participant outcome.

The candidate is not yet a publicly publishable production release. Production signing, stable public policy and support endpoints, console declarations, internal-track evaluation, and store review remain external owner actions. This distinction is part of the result rather than an administrative footnote: trustworthy privacy software must disclose where its evidence ends.

\subsection{Implications for Privacy-Engineering Practice}

The first lesson is that provenance must survive every transformation. Capturing a source label at ingestion is insufficient if storage, notifications, exports, or screenshots omit it. Provenance should be treated as part of the value being computed. Tests should therefore follow it from observation through decision and presentation.

The second lesson is that missing context requires a first-class state. Many privacy questions cannot be answered from technical frequency. Returning a neutral or positive result when purpose, authority, or acquisition coverage is unknown creates hidden assumptions. An evidence-gap state makes the limitation actionable: it can ask for a privacy notice, controller record, consent basis, or authorised trace rather than fabricating certainty.

The third lesson is that extensibility needs governance. An interface or schema that can load new legal text is technically flexible but may accelerate error if sources, versions, review, and updates are uncontrolled. The pack boundary is valuable precisely because it creates a unit that can be admitted, tested, rejected, and deprecated.

The fourth lesson is that polished design should reduce ambiguity rather than conceal it. Visual hierarchy can foreground evidence strength, state, and next action. Calm colour and concise language can lower cognitive burden. Yet critical caveats must remain visible. The ethical objective is not to maximise engagement or confidence; it is to help the user reach an appropriately bounded understanding.

The fifth lesson is that store readiness is an evidence package. An AAB is necessary but not sufficient. Signing ownership, manifest reconciliation, data declarations, policy endpoints, screenshots, accessibility, update testing, and support processes all affect whether the distributed product matches the evaluated one. Treating these as afterthoughts would break the thesis's claim chain at the point of deployment.

The final lesson is that disciplined limitation can increase, rather than reduce, practical value. A reviewer can act on a clearly sourced question, ask for missing evidence, and revisit the decision after a pack update. An unsupported compliance badge offers none of those affordances. Privacy Lens is most credible when it remains specific about what it observed, how it reasoned, and what a human must still decide.

\subsection{Future Work}

Future work should commission qualified review of the GDPR pack and its source classifications, automate monitored retrieval and hashing from official authorities, define a signed approval and deprecation workflow, add a genuinely reviewed regional pack, improve authorised event-level acquisition, test multiple Android implementations over long durations, and conduct accessibility and participant studies. These activities should retain the Revision 19 claim-to-evidence discipline and its separation between implemented governance, legal validity, and observed human outcomes. Accessibility work should add TalkBack, switch access, magnification, localisation, orientation, and assistive-technology participants. A participant protocol should test signal-versus-verdict comprehension, recognition of missing context, source-finality comprehension, proportional action choice, calibrated reliance, task completion, and willingness to continue.

The programme should also investigate formal and property-based techniques for evidence admission and replay. Useful properties include idempotence, bounded storage, monotonic watermarks, source isolation, and stability under reordered independent windows. Mutation testing could assess whether the suite detects weakened validation. A separately implemented oracle would reduce the correlated-defect risk of deriving implementation and expected values from one specification.

Interface research should compare alternative presentations of uncertainty. For example, a three-state label, an evidence-strength scale, and a question-led presentation may produce different reliance patterns. The comparison should measure comprehension and decisions, not only preference. Accessibility evaluation should be incorporated from recruitment through reporting rather than added as a final conformance scan.

The pack framework could support signed declarative updates after its threat model is tested. This would require publisher identity, schema versioning, offline recovery, rollback protection, and clear user communication. Remote executable rules would introduce unnecessary risk; the safest extension preserves a narrow validated data contract.

Finally, future scholarship should study the relationship between accountability prompts and organisational action. A technically clear question has value only if users or reviewers can obtain the missing policy, consent, retention, or controller evidence. Research with data-protection practitioners could examine which prompts accelerate review, which create noise, and how findings should integrate with established assessment records without becoming an unreviewed legal decision.

\subsection{Closing Statement}

The central finding is that privacy accountability on Android is strengthened not by presenting the strongest possible conclusion, but by presenting the strongest conclusion the available evidence actually supports. Privacy Lens operationalises that principle in code, interface, evaluation, and thesis structure.
